% ============================================================================
\section{验证}
\label{sec:tss-verify}
% ============================================================================

我们通过 \texttt{ring3 test} shell 命令来验证整个 TSS + Ring 3 机制。

% ----------------------------------------------------------------------------
\subsection{测试流程}
% ----------------------------------------------------------------------------

\codefile{src/kernel/cmds/cmd\_ring3.c}
\begin{lstlisting}[style=cstyle]
static int ring3_test(void) {
    /* 1. 分配物理页并映射到用户代码地址 */
    uint32_t code_phys = (uint32_t)(uintptr_t)pmm_alloc_page();
    vmm_map_page(USER_CODE_VADDR, code_phys,
                 PTE_PRESENT | PTE_USER);

    /* 2. 分配物理页并映射到用户栈地址 */
    uint32_t stack_phys = (uint32_t)(uintptr_t)pmm_alloc_page();
    vmm_map_page(USER_STACK_VADDR, stack_phys,
                 PTE_PRESENT | PTE_WRITABLE | PTE_USER);

    /* 3. 将 HLT 指令（0xF4）复制到用户代码页 */
    kmemcpy((void*)USER_CODE_VADDR,
            user_test_code, sizeof(user_test_code));

    /* 4. 注册 VMA（便于 #PF handler 诊断） */
    vma_add(USER_CODE_VADDR, ..., VMA_READ | VMA_EXEC, "user-code");
    vma_add(USER_STACK_VADDR, ..., VMA_READ | VMA_WRITE, "user-stack");

    /* 5. 跳转到 Ring 3 —— 不返回！ */
    jump_to_ring3(USER_CODE_VADDR, USER_STACK_TOP);
}
\end{lstlisting}

地址选择：

\begin{itemize}
  \item 用户代码页：\hex{00400000}（4MiB 处，避开 NULL 区域）
  \item 用户栈页：\hex{00401000}，栈顶 = \hex{00402000}（栈向低地址增长）
\end{itemize}

% ----------------------------------------------------------------------------
\subsection{用户态测试代码}
% ----------------------------------------------------------------------------

测试代码只有 3 字节：

\begin{lstlisting}[style=nasmstyle]
0xF4        ; hlt — Ring 3 执行此指令触发 #GP
0xEB, 0xFE  ; jmp $ — 无穷循环（安全网）
\end{lstlisting}

\noindent \texttt{HLT} 是特权指令——只有 Ring 0 可以执行。
Ring 3 执行 \texttt{HLT} 时，CPU 会触发 \#GP（General Protection Fault, vector 13）。

% ----------------------------------------------------------------------------
\subsection{预期输出}
% ----------------------------------------------------------------------------

\begin{tcblisting}{consolebox, title={\tpk{>}~ring3 test}}
«\tpk{szy-kernel >}» «\tin{ring3 test}»
[ring3] Preparing Ring 3 test...
[ring3] code page: virt=0x00400000 phys=0x001b1000
[ring3] stack page: virt=0x00401000 phys=0x001b3000 (top=0x00402000)
[ring3] test code copied (3 bytes): HLT instruction
[ring3] Jumping to Ring 3 (eip=0x00400000, esp=0x00402000)...
[ring3] Expected: #GP fault (HLT is privileged)

[isr] #GP General Protection vector=13 err=0x00000000 eip=0x00400000
[isr] Fault from Ring 3 (user mode)
[isr] User CS=0x001b  SS=0x0023  ESP=0x00402000
Kernel panic: #GP General Protection
\end{tcblisting}

关键验证要点：

\begin{enumerate}
  \item \texttt{Fault from Ring 3 (user mode)} —— ISR 正确识别异常来自用户态
  \item \texttt{CS=0x001b} —— 用户代码段选择子（GDT[3] | RPL=3）
  \item \texttt{SS=0x0023} —— 用户数据段选择子（GDT[4] | RPL=3）
  \item \texttt{eip=0x00400000} —— 故障发生在用户代码页的入口地址
  \item \texttt{err=0x00000000} —— \#GP 错误码为 0 表示常规保护违规
\end{enumerate}

\begin{whybox}
为什么 ISR 能正确打印"Fault from Ring 3"？

当 \#GP 从 Ring 3 触发时，CPU 自动执行以下操作：
\begin{enumerate}
  \item 从 TSS 读取 SS0:ESP0，切换到内核栈
  \item 在内核栈上压入旧的 SS/ESP/EFLAGS/CS/EIP 和 error code
  \item 跳转到 IDT[13] 的处理程序（我们的 \texttt{isr13}）
\end{enumerate}

ISR 的 \texttt{isr\_common\_stub} 把这些值打包成 \texttt{regs\_t} 结构体传给 C。
C 代码检查 \texttt{r->cs \& 0x3}：CS 的低 2 位 = RPL = 异常发生时的 CPL。
如果 RPL=3，就打印"Ring 3"。

结构体末尾的 \texttt{user\_esp} 和 \texttt{user\_ss} 字段在跨特权级中断时
自动出现在栈上（CPU 额外压入），同特权级时不存在。
\end{whybox}
